%قبل از شروع مسئله ترجیحاً چیزی ننویسید. اما اگر خواستید بنویسید، در قالب لیستی (فقرات و شمارش) بنویسید.
%برای اضافه کردن نمره، بعدی از شروع‌مسئله، [x] را اضافه کنید و به جای x  نمره مورد نظر را قرار دهید. برای زیرمسئله‌ها هم به همین صورت است.
%\شروع‌مسئله[10]
\شروع‌مسئله [26]
% صورت اصلی مسئله را در این‌جا بنویسید.
	

%اگر مسئله چند قسمت داشت به این صورت عمل کنید:
\شروع{زیرمسائل}
\زیرمسئله[16] % مسئلۀ ۱آ
یک درخت دودویی جستجوی کامل داریم. الگوریتمی با زمان اجرای
$\mathcal{O}(n)$
طراحی کنید که این درخت دودویی جستجو را به هرم بیشینه‌ای تبدیل کند که گره‌های زیردرخت راست هر رأس آن از گره‌های زیردرخت سمت چپش بزرگ‌تر باشند (در هر گام می توانیم جای دو رأس را با هم عوض کنیم).

\پاسخ{
واضح است که پیمایش پس‌ترتیب هرم بیشینه عناصر را به صورت مرتب شده به ما تحویل می‌دهد. در مورد درخت دودویی جستجو هم می‌دانیم که پیمایش میان‌ترتیب آن عناصر را به صورت مرتب‌شده به ما می‌دهد. پس پیمایش میان‌ترتیب د.د.ج با پیمایش پس‌ترتیب هرم یکسان است. در نتیجه می‌توانیم مقدار هر رأس هرم بیشینه را در
$O(n)$
به‌دست آوریم (چون پیمایش یک درخت در زمان 
$O(n)$
انجام می شود)
و بعد از آن هم به راحتی با عوض کردن هر رأس با رأس مورد نظر مسئله حل می‌شود.
	
}

\زیرمسئله[10] % مسئلۀ ۱ب
ثابت کنید برای تبدیل یک درخت دودویی متوازن که درخت آن تقریباً کامل نیز هست (مثل هرم) به هرم بیشینه‌ای مشابه قسمت آ، الگوریتمی بهتر از
$\mathcal{O}(n\log n)$
وجود ندارد (مانند قسمت قبل در هر مرحله تنها می‌توانیم جای دو رأس را با هم عوض کنیم).

\پاسخ{
	پیمایش پس ترتیب هرم را در نظر بگیرید و آن را $A$ بنامید. عوض کردن جای دو رأس مثل این است که جای دو عدد را در $A$ تغییر بدهیم. برای این‌که شرط مسئله محقق شود باید $A$ مرتب شود در حالی که در ابتدا اعداد موجود در آرایۀ $A$ هیچ شرط خاصی ندارند. می‌دانیم این کار در زمانی کمتر از 
	$O(n\log n)$
	میسر نیست.
	
}
%\زیرمسئله % مسئلۀ ۱پ
%\زیرمسئله % مسئلۀ ۱ت
\پایان{زیرمسائل}





%%%%%%%%%%%%%%%%%%%%%%%%%%%%%%%%%%%%%%%%%%%%%%%%%%%%%
% برخی دستورات مفید برای تمرین‌ها. %
%%%%%%%%%%%%%%%%%%%%%%%%%%%%%%%%%%%%%%%%%%%%%%%%%%%%%

% فرمول‌های نمونه:
%$f(x) = \Theta(x)$
%$T(x, y) \leq \log(x) + 2^y + \binom{2n}{n}$
% {\tt code\_function}


% لیست موارد بدون شماره‌گذاری:
%\شروع{فقرات}
%    \فقره مورد اول 
%        \شروع{فقرات}
%            \فقره مورد اول داخلی 
%                \شروع{فقرات}
%                    \فقره اولین مورد خیلی داخلی 
%                    \فقره دومین مورد خیلی داخلی 
%                \پایان{فقرات}
%            \فقره مورد دوم داخلی 
%        \پایان{فقرات}
%    \فقره مورد دوم 
%\پایان{فقرات}

% لیست موارد شماره‌گذاری شده:
%\شروع{شمارش}
%    \فقره یک مورد 
%    \فقره یک مورد دیگر 
%    \فقره باز هم یک مورد دیگر
%\پایان{شمارش}

% برای نوشتن یک فرمول چند خطی که یک شماره بخوره. برای حالت‌های دیگر به سایت زیر مراجعه کنید:
% https://www.overleaf.com/learn/latex/Aligning_equations_with_amsmath
%\begin{align} 
%    \begin{split}
%        (x+y)^3 &= (x+y)^2(x+y) \\
%                &= (x^2+2xy+y^2)(x+y) \\
%                &= (x^3+2x^2y+xy^2) + (x^2y+2xy^2+y^3) \\
%                &= x^3+3x^2y+3xy^2+y^3
%    \end{split}                                 
%\end{align}

% نوشتن ماتریس و آرایه دوبعدی
%\begin{align}
%    A = 
%    \begin{bmatrix}
%        A_{11} & A_{21} \\
%        A_{21} & A_{22}
%    \end{bmatrix}
%\end{align}

% اضافه کردن تصویر با شرح. برای اضافه کردن تصویر بدون شرح، فقط دستور درج‌تصویر را استفاده کنید.
%\شروع{شکل}[h!]
%\تنظیم‌ازوسط
%  \درج‌تصویر[پهنا=0.5\پهنای‌متن]{img.jpg}
%\شرح{شرح تصویر}
%\برچسب{تصویر:۱}
%\پایان{شکل}
%\پایان{زیرمسائل}

%کد زیر را ببینید.
%%\lstinputlisting{source_filename.py}
%\begin{latin}
%\begin{lstlisting}
%# Import the modules
%import sys
%import random
%
%ans = True
%
%while ans:
%question = raw_input("Ask the magic 8 ball a question: (press enter to quit) ")
%
%answers = random.randint(1,8)
%
%if question == "":
%sys.exit()
%
%elif answers == 1:
%print "It is certain"
%
%elif answers == 2:
%print "Outlook good"
%
%elif answers == 3:
%print "You may rely on it"
%
%elif answers == 4:
%print "Ask again later"
%
%elif answers == 5:
%print "Concentrate and ask again"
%
%elif answers == 6:
%print "Reply hazy, try again"
%
%elif answers == 7:
%print "My reply is no"
%
%elif answers == 8:
%print "My sources say no"
%\end{lstlisting}
%\end{latin}
